\section{Technical Design}
To ensure that effort expended on this project is efficient,  it's important
to align the project's objectives/goals, with its design and therefore the tools that shall fulfil that design.
Thus, here we discuss each objective and its derivative design along with the technology that will
achieve said goal (e.g. programming language, framework, library etc.). 
Docker

Table 1 showcases this.

\begin{table}[h!]
\begin{center}
\begin{tabular}{||c | c | c||}
    \hline
    Goal & Design & Tool \\

    \hline
    Multiplayer & Authoritative Server \& Database & NoSQL \\
    \hline
    AI          & Backend Minimax \& MCTS & Golang/Rust Backend \\
    \hline
    Website features & Typescript \& Web frameworks & Dart/Flutter \\ 
    \hline 
\end{tabular}
\end{center}
\caption{Design overview.}
\label{table:1}
\end{table}

Next, we break down each goal, the design decision and the respective tool. 

\subsection {Multiplayer: Authoritative Server \& Database}
To achieve multiplayer, there must be a way to synchronise multiple instances of games, and player actions with each other.
Furthermore, an uncomprimisable tenet of multiplayer for player versus player games is competitive integerity. Thus,
a server model which ensures said integrity is paramount to the project's success.

Therefore, an authoritative server model in which a central trusted node cross-checks and resolves game state differences, unintentionally or intentionally through cheating,
is best to fulfil the promise and need for competitive integrity. 

\subsection {AI: Minimax \& MCTS Server Backend}
This section briefly explores the API's shape. 

\subsection{Architecture}
Here we explore how the \verb|ut3-oxide| Rust implementation will operate.

\subsubsection{Binary Operations}
For maximal efficiency's sake, we opt for representing each row as single 
\verb|u32| integer, with 2 bits per cell ($9 \times 2 = 18$ — leaving 14 bits remaining for metadata)
\paragraph{Getting Bits}
\begin{center}
    \bitpattern[numberBitsBelow,numberFieldsTwice,startBit=31]{meta data}[14]{cells}[18]/
\end{center}

Formula
\begin{center}
    \begin{math}
        n \& ((2^x - 1) \ll dx) \gg dx
    \end{math}
\end{center}

Alternatively
\begin{center}
\begin{math}
    n \operatorname{\&} (((1 \ll x) - 1)) \ll dx \gg dx
\end{math}
\end{center}

NB: Naturally, if you don't need to offset the value, then $dx$ is unrequired and is simplified to:

\begin{center}
    \begin{math}
        n \operatorname{\&} (1 \ll x) - 1
    \end{math}
\end{center}

\paragraph{Clearing Bits}
Here's how bits are cleared
\begin{center}
    $n = 1101 1100$
    

    $mask = 0b11 \ll dx$
    
    
    $n = n \operatorname{\&} \operatorname{\sim} mask$
\end{center}

\begin{center}
    \begin{math}
        n = n \land \neg(0b11 \ll dx)
    \end{math}
\end{center}

    n \&= $\sim (0b11 \ll dx)$
\paragraph{Setting Bits}
\subsection{Components}
\subsection{Algorithms}
As previously mentioned, there will be two AI Algorithms for the Player vs Computer mode.
Furthermore, it's the Minimax Algorithm that ought ot be used for the easy difficulty. However,
the relative difficulty depends on minimax's evaluation function, and the duration allowed for 
MCTS to explore the game tree. Therefore, it may be more efficient to develop the MCTS AI more thoroughly and create 
a perceived difficult through differing allowed tree-search time durations. 

\subsubsection{MCTS}

\subsection{Data Models}
\subsection{APIs and Interfaces}

\subsection{Tech Stack}
Frontend: Flutter (Dart) \newline
Backend: Rust SqlX?
